%%
% 第三章：实验过程与结果分析
%%

\chapter{实验过程与结果分析}
\label{chap:experiment}

\section{实验设计与评价指标}

为保证实验结果具有可重复性，本章采用\textbf{受控实验}方式对 \path{agent_pageindex.py} 进行覆盖率分析。被测对象仍然是 PageIndex 的真实入口脚本，但实验时不直接依赖外部网络调用，而是通过桩替换的方式固定 Agent 和工具返回值。这样处理的目的，是把实验变量收紧，让每一组用例都对应一个明确的控制流目标。

本章实验主要回答三个问题：

\begin{itemize}
    \item PathTracer 能否稳定识别 \path{agent_pageindex.py} 中全部 22 个控制结构入口；
    \item 在普通模式、流式模式和列表型消息输入三种场景下，脚本分别会走到哪些路径；
    \item 三组测试用例的结果并集能否覆盖全部 22 个控制结构入口。
\end{itemize}

本章继续使用第二章定义过的控制结构命中率作为核心指标：

\[
\text{Coverage} = \frac{\text{ExecutedBranches}}{\text{TotalBranches}} \times 100\%
\]

这个公式里的分子表示“控制结构起始行被执行到”的数量，%
分母表示静态分析阶段识别出的控制结构总数。%
这个口径依然不是严格意义上的真假分支覆盖，因此后文即使出现 100\% 覆盖率，%
也只能说明 22 个控制结构入口都被执行到了，不能说明每个判断的真分支和假分支都被分别验证过。

\section{实验环境与实验对象}

\subsection{实验环境}

实验在当前项目环境中完成，关键配置如表 \ref{tab:exp-env} 所示。

\begin{table}[htbp]
    \centering
    \small
    \caption{实验环境配置}
    \label{tab:exp-env}
    \begin{tabular}{p{0.23\textwidth}p{0.62\textwidth}}
        \toprule
        项目 & 配置 \\
        \midrule
        运行平台 & Darwin arm64 \\
        Python 版本 & 3.13.11 \\
        主要依赖 & \path{langchain>=1,<2}，\path{langchain-openai}，\path{python-dotenv==1.1.0} \\
        被测脚本 & \path{agent_pageindex.py} \\
        追踪工具 & 自研工具 \texttt{PathTracer}（包名 \path{pathtracer}） \\
        实验驱动脚本 & \path{tests/run_agent_pageindex_experiments.py} \\
        结果文件 & \path{docs/coverage_tc1.json}、\path{docs/coverage_tc2.json}、\path{docs/coverage_tc3.json} \\
        \bottomrule
    \end{tabular}
\end{table}

其中，\path{tests/run_agent_pageindex_experiments.py} 是本章实验的驱动脚本。它不会修改 \path{agent_pageindex.py} 本身，而是在实验进程里临时安装假的 \path{langchain}、\path{dotenv} 和 \path{pageindex.search} 模块，再通过 \path{runpy.run_path(..., run_name="__main__")} 的方式直接执行真实脚本。这样既保住了“被测对象是 PageIndex 入口脚本”这一前提，又把外部返回值变成了可精确控制的输入。

\subsection{被测对象说明}

\path{agent_pageindex.py} 的总代码量约 130 行，控制结构主要集中在三个函数中：

\begin{itemize}
    \item \path{_extract_text}：负责统一处理 \path{None}、字符串、列表等不同类型的消息内容；
    \item \path{_print_verbose_stream}：负责消费流式消息，内部包含多层 \texttt{for} 循环、重复消息过滤、工具消息和模型消息分流；
    \item \path{main}：负责参数解析、环境变量检查，以及普通模式和 \path{verbose} 模式之间的主流程切换。
\end{itemize}

静态分析结果表明，该脚本共包含 22 个可统计的控制结构入口，其中 \texttt{if} 结构 17 个，\texttt{for} 结构 5 个。按函数拆分后的分布如表 \ref{tab:static-branch-distribution} 所示。

\begin{table}[htbp]
    \centering
    \caption{agent\_pageindex.py 静态控制结构分布}
    \label{tab:static-branch-distribution}
    \begin{tabular}{lccc}
        \toprule
        函数 & \texttt{if} 入口数 & \texttt{for} 入口数 & 合计 \\
        \midrule
        \path{_extract_text} & 5 & 1 & 6 \\
        \path{_print_verbose_stream} & 7 & 4 & 11 \\
        \path{main} & 5 & 0 & 5 \\
        \midrule
        总计 & 17 & 5 & 22 \\
        \bottomrule
    \end{tabular}
\end{table}

这个分布说明，\path{_print_verbose_stream} 和 \path{_extract_text} 两部分集中了脚本中大多数控制结构。如果消息形态不可控，这两部分就很容易留下明显覆盖空洞。

\section{测试用例设计}

\subsection{设计原则}

本章测试用例围绕\textbf{目标路径}进行设计。每组用例只负责覆盖一类关键控制结构，三组结果求并集后再评估 22 个控制结构入口的整体覆盖情况。

设计原则如下：

\begin{itemize}
    \item 被测对象固定为 \path{agent_pageindex.py}，不换成 demo 脚本；
    \item 外部依赖全部桩化，避免网络和 LLM 输出随机性影响结果；
    \item 每组用例只改变一种关键消息形态，便于解释覆盖率增量来自哪里；
    \item 所有结果统一输出为 JSON，方便直接做行号集合分析。
\end{itemize}

\subsection{测试用例定义}

三组实验场景如表 \ref{tab:testcases} 所示。

\begin{table}[htbp]
    \centering
    \small
    \caption{实验测试用例}
    \label{tab:testcases}
    \begin{tabular}{p{0.12\textwidth}p{0.26\textwidth}p{0.44\textwidth}}
        \toprule
        用例 & 场景名 & 设计目的 \\
        \midrule
        TC1 & \path{normal\_string} & 普通模式下返回字符串型最终答案，作为普通调用路径基线 \\
        TC2 & \path{verbose\_mixed} & 流式模式下返回混合消息，覆盖非字典块、重复消息、工具调用、空消息、\path{ToolMessage} 和 \path{AIMessage} 路径 \\
        TC3 & \path{normal\_list} & 普通模式下返回列表型消息内容，专门覆盖 \path{_extract_text} 的列表输入路径 \\
        \bottomrule
    \end{tabular}
\end{table}

这里最关键的是 TC3。该用例专门针对 \path{_extract_text} 的列表输入路径设计，使最终消息内容固定为：

\begin{lstlisting}[language={} ,caption={TC3 中固定的列表型消息内容}]
["alpha fragment", {"type": "text", "text": "beta fragment"}]
\end{lstlisting}

这样一来，\path{_extract_text} 中“判断为列表”“遍历列表”“字符串元素处理”“字典元素处理”四个控制结构都能稳定触发。

\section{实验执行过程}

\subsection{执行命令}

三组实验通过同一个驱动脚本执行，对应命令如下：

\begin{lstlisting}[language=bash, caption={实验执行命令}]
python3 tests/run_agent_pageindex_experiments.py \
    --scenario normal_string --output docs/coverage_tc1.json

python3 tests/run_agent_pageindex_experiments.py \
    --scenario verbose_mixed --output docs/coverage_tc2.json

python3 tests/run_agent_pageindex_experiments.py \
    --scenario normal_list --output docs/coverage_tc3.json
\end{lstlisting}

\subsection{执行步骤}

每组实验都按相同步骤完成：

\begin{enumerate}
    \item 在实验进程中安装假的 \path{langchain}、\path{dotenv} 和 \path{pageindex.search} 模块；
    \item 根据场景名创建对应的假 Agent；
    \item 用 PathTracer API 将追踪目标固定为 \path{agent_pageindex.py}；
    \item 通过 \path{runpy.run_path} 以 \path{__main__} 方式执行目标脚本；
    \item 用 \path{CoverageReporter} 对执行结果做静态-动态对齐，并输出 JSON 报告。
\end{enumerate}

这里值得强调一点：实验虽然用了驱动脚本，但真正被追踪的文件始终是 \path{agent_pageindex.py}，而不是驱动脚本本身。因此，报告中的行号、分支数和覆盖率都仍然属于 PageIndex 入口脚本。

\section{实验结果}

\subsection{总体覆盖率结果}

三组实验结果如表 \ref{tab:pageindex-coverage} 所示。

\begin{table}[htbp]
    \centering
    \caption{PageIndex 实验覆盖率结果}
    \label{tab:pageindex-coverage}
    \begin{tabular}{lcccc}
        \toprule
        测试用例 & 已命中分支数 & 总分支数 & 覆盖率 & 执行行数 \\
        \midrule
        TC1 & 7 & 22 & 31.8\% & 55 \\
        TC2 & 16 & 22 & 72.7\% & 81 \\
        TC3 & 11 & 22 & 50.0\% & 62 \\
        \bottomrule
    \end{tabular}
\end{table}

TC1 仍然是最低的，只覆盖了普通模式下的基础路径：参数解析、环境检查、\path{agent.invoke} 返回结果读取，以及 \path{_extract_text} 的字符串输入处理。它的作用更多是作为基线，而不是承担补覆盖的任务。

TC2 是增量最大的用例，覆盖率达到 72.7\%。这一组新增的 16 个命中分支由事先设计好的混合流式消息稳定触发。具体来说，这一组消息同时覆盖了：

\begin{itemize}
    \item \path{chunk} 不是字典时的过滤；
    \item \path{payload} 不是字典时的过滤；
    \item 重复消息的去重；
    \item 工具调用消息与工具调用列表遍历；
    \item 空内容消息的跳过；
    \item \path{ToolMessage} 和 \path{AIMessage} 两类消息的分流打印。
\end{itemize}

TC3 的覆盖率是 50.0\%，虽然低于 TC2，但它补的是 \path{_extract_text} 的四个关键列表路径分支。对整个实验设计来说，TC3 的价值不在于单次覆盖率高，而在于它把普通字符串输入难以覆盖的列表路径单独固定下来，形成可重复命中的测试对象。

\subsection{分支交集与并集分析}

为了更清楚地看到每组用例各自贡献了什么，把三组结果的分支集合拆开后可得到表 \ref{tab:branch-set-relation}。

\begin{table}[htbp]
    \centering
    \small
    \caption{三组测试用例的分支集合关系}
    \label{tab:branch-set-relation}
    \begin{tabular}{p{0.18\textwidth}p{0.25\textwidth}p{0.44\textwidth}}
        \toprule
        类别 & 行号 & 对应逻辑 \\
        \midrule
        三组共有 & 25, 27, 88, 116, 129 & 文本提取入口、字符串判定、API Key 检查、\path{verbose} 判定和脚本入口判断 \\
        仅 TC2 命中 & 43, 44, 46, 47, 50, 52, 57, 58, 64, 67, 73 & 流式块遍历、非法块过滤、消息循环、去重、工具调用、空消息处理、工具消息和模型消息分流 \\
        仅 TC3 命中 & 29, 31, 32, 34 & \path{_extract_text} 的列表判定、列表遍历、字符串元素处理和字典元素处理 \\
        仅 TC1 命中 & 无 & 基线用例本身不提供独占增量 \\
        \bottomrule
    \end{tabular}
\end{table}

从这个结果可以看出，TC1 主要提供普通模式基线，本身不带来独占增量；决定整体覆盖效果的是 TC2 和 TC3。TC2 负责覆盖 \path{_print_verbose_stream} 中的流式处理逻辑，TC3 负责覆盖 \path{_extract_text} 的列表路径，两者分工清晰，也符合“围绕控制结构设计实验”的思路。

\subsection{累计覆盖结果}

把三组实验的已命中分支做并集统计之后，可以得到完整的 22 个控制结构入口：

\begin{lstlisting}[language={} ,caption={三组实验结果的分支并集}]
25, 27, 29, 31, 32, 34, 43, 44, 46, 47, 50, 52,
57, 58, 64, 67, 73, 88, 116, 121, 123, 129
\end{lstlisting}

这意味着三组实验结果求并集后，累计覆盖率达到：

\[
\frac{22}{22} \times 100\% = 100.0\%
\]

这里的“100\%”只对应控制结构入口的全覆盖。它说明 \path{agent_pageindex.py} 中所有 22 个静态控制结构都至少被执行到过一次；但它不意味着真假分支方向都被完整区分，也不意味着底层 PageIndex 检索模块已经被完整测透。这一点后面还会再说明。

\section{结果解释与路径分析}

\subsection{TC1：普通模式基线路径}

TC1 的路径最简单。程序进入 \path{main} 后，完成参数解析、环境变量检查和 Agent 创建，随后走普通模式的 \path{agent.invoke} 分支，在第 121 行和第 123 行判断返回结果是否为字典、是否包含消息列表，最后把最后一条消息内容送进 \path{_extract_text}。由于返回内容为普通字符串，所以 \path{_extract_text} 只会停留在第 25 行和第 27 行附近，不会进入列表路径。

这个用例的意义在于给后两组实验提供对照。它告诉我们：如果只验证最常见的普通调用路径，入口脚本里其实还有一大半结构不会被触发。

\subsection{TC2：流式混合消息路径}

TC2 是本组实验的主力用例。它把程序显式带进 \path{verbose} 模式，并且在流式返回里依次制造了五类事件：

\begin{itemize}
    \item 非字典 \path{chunk}，触发第 44 行过滤；
    \item 字典 \path{chunk} 但 \path{payload} 不是字典，触发第 47 行过滤；
    \item 带 \path{tool_calls} 的消息，触发第 57 行和第 58 行；
    \item 重复消息，触发第 52 行；
    \item 空内容消息、\path{ToolMessage} 和 \path{AIMessage}，分别触发第 64、67、73 行。
\end{itemize}

这种设计比单纯加一个 \path{--verbose} 参数更有效，因为它不是笼统地进入流式模式，而是把流式模式内部的关键分支一个个拆开打。这样得到的 72.7\% 覆盖率，不依赖真实消息流的随机形态，而是直接由测试数据控制。

\subsection{TC3：列表型消息内容路径}

TC3 是本章实验中用于补齐列表处理路径的关键用例。它依然走普通模式，但把最终消息内容固定成“字符串 + 文本字典”的列表。这样程序在调用 \path{_extract_text} 时，会连续经过以下结构：

\begin{itemize}
    \item 第 29 行：判断 \path{payload} 是 \path{list}；
    \item 第 31 行：遍历列表元素；
    \item 第 32 行：处理字符串元素；
    \item 第 34 行：处理 \path{\{"type": "text"\}} 形式的字典元素。
\end{itemize}

这四处结构都集中在同一条列表输入路径上。TC3 把这条路径单独抽出来验证，所以它虽然单次覆盖率只有 50.0\%，但在整个实验设计里反而是不可替代的一组。

\subsection{实验设计说明}

本组实验之所以采用受控输入而不是直接依赖真实查询返回，主要有两个原因。

第一，变量更可控。查询参数、外部返回形态和脚本内部路径这三者如果同时变化，覆盖率差异就很难解释清楚；将外部返回固定为预设场景后，每个增量都可以明确对应到具体控制结构。

第二，结论更可复查。三组结果都可以通过 \path{tests/run_agent_pageindex_experiments.py} 重新生成，不依赖网络、不依赖模型当时的随机输出，也不依赖某次数据库返回恰好是什么内容。对于课程实验报告来说，这一点非常重要，因为它让“为什么打到了这些行”变成了可验证的事实，而不是一次偶然运行的现象。

\section{实验结论与局限}

\subsection{实验结论}

本章实验表明，PathTracer 可以在 PageIndex 真实入口脚本 \path{agent_pageindex.py} 上稳定工作，并通过受控输入把 22 个控制结构入口全部命中。具体来说：

\begin{itemize}
    \item TC1 覆盖普通模式基线路径，覆盖率为 31.8\%；
    \item TC2 覆盖流式处理路径，覆盖率为 72.7\%；
    \item TC3 覆盖列表型消息内容路径，覆盖率为 50.0\%；
    \item 三组结果取并集后，累计覆盖率达到 100.0\%。
\end{itemize}

这说明 PathTracer 至少在入口脚本这一层，已经具备了较完整的控制结构检测能力。更重要的是，这一覆盖结果来自明确的测试目标与受控输入设计，而不是依赖外部返回内容的偶然形态。

\subsection{仍然存在的局限}

即使本章已经拿到了 100.0\% 的控制结构命中率，也不能把这个结果说得过头。它至少还有两个边界。

\begin{itemize}
    \item 第一，这个 100.0\% 只代表控制结构入口全覆盖，不代表真假分支方向全覆盖。第 88 行和第 116 行就是最典型的例子：这一行判断被执行到了，但条件真假并没有被分别统计。
    \item 第二，本章重点覆盖的是 \path{agent_pageindex.py} 入口脚本，而不是 \path{pageindex/search.py} 这类更深层模块。也就是说，当前实验主要验证了“入口脚本路径是否可测”，还没有展开到底层检索逻辑的更深层覆盖分析。
\end{itemize}

这两个限制并不削弱本章的主要结论，但需要如实写出来。否则只报一个 100\%，很容易让人误以为整个 PageIndex 项目都已经被充分测试了。

\section{本章小结}

本章围绕 PageIndex 入口脚本 \path{agent_pageindex.py} 设计了一组可重复的受控实验方案。该方案保持被测对象不变，但把外部依赖全部桩化，并通过 \path{tests/run_agent_pageindex_experiments.py} 生成可复现的覆盖率结果。三组用例分别覆盖普通模式基线路径、流式混合消息路径和列表型消息内容路径，累计覆盖率达到 100.0\%。

更关键的是，本章把“覆盖率为什么会变化”讲清楚了：TC2 负责打通流式处理逻辑，TC3 负责补齐 \path{_extract_text} 的列表路径，二者合起来才真正把 22 个控制结构入口全部覆盖。这样得到的实验结果更适合作为正式实验章节的依据。
